\documentclass[14pt, a4paper, oneside]{extreport}
\usepackage{extsizes}

\usepackage{cmap} % Улучшает отображение символов в PDF
% \usepackage{plantuml}

% \usepackage{microtype}       % Улучшает переносы и плотность текста
% \usepackage[htt]{hyphenat}   % Позволяет переносить слова с дефисами и URL

\usepackage[T2A]{fontenc}
\usepackage[utf8]{inputenc}
\usepackage[russian]{babel}
\usepackage{geometry}
\geometry{left=30mm, right=10mm, top=20mm, bottom=20mm}
\usepackage{setspace}
\onehalfspacing
\setlength{\parindent}{1.25cm}
\setlength{\parskip}{0pt}
\usepackage{indentfirst}

% ==========================================
% ФОРМАТИРОВАНИЕ (ВСЁ В ПРЕАМБУЛЕ)
% ==========================================
\tolerance=1000
\emergencystretch=1em
\hbadness=10000

\usepackage{tabularx}
\usepackage{amsmath, amssymb, mathtools}
% \numberwithin{equation}{section} % Нумерация формул (1.1), (2.3)

\usepackage{titlesec}
\titleformat{\chapter}[block]{\normalfont\normalsize\normalsize}{}{0pt}{\hspace{1.25cm}}
\titleformat{\section}[block]{\normalfont\normalsize\normalsize}{}{0pt}{\hspace{1.25cm}}
\titleformat{\subsection}[block]{\normalfont\normalsize\normalsize}{}{0pt}{\hspace{1.25cm}}
\titlespacing*{\chapter}{0pt}{0pt}{1em}
\titlespacing*{\section}{0pt}{1em}{0.5em}
\titlespacing*{\subsection}{0pt}{1em}{0.5em}

\usepackage{tocloft}

% Название оглавления (правильный способ для babel)
\addto\captionsrussian{\renewcommand{\contentsname}{СОДЕРЖАНИЕ}}
\addto\captionsrussian{\renewcommand{\bibname}{СПИСОК ИСПОЛЬЗОВАННЫХ ИСТОЧНИКОВ}}

% Шрифт заголовка "СОДЕРЖАНИЕ" (строго 14pt, жирный, с красной строкой)
\renewcommand{\cfttoctitlefont}{\normalfont\normalsize\hspace{1.25cm}}
\setlength{\cftbeforetoctitleskip}{0pt}
\setlength{\cftaftertoctitleskip}{1em}

% Отступы для записей оглавления (красная строка 1.25 см)
\setlength{\cftchapindent}{0cm}
\setlength{\cftsecindent}{0cm}
\setlength{\cftsubsecindent}{0cm}

% Ширина поля под номер (0pt, так как нумерация ручная)
\setlength{\cftchapnumwidth}{0pt}
\setlength{\cftsecnumwidth}{0pt}
\setlength{\cftsubsecnumwidth}{0pt}

% Точки-заполнители до номера страницы
\renewcommand{\cftchapleader}{\cftdotfill{\cftdotsep}}
\renewcommand{\cftsecleader}{\cftdotfill{\cftdotsep}}
\renewcommand{\cftsubsecleader}{\cftdotfill{\cftdotsep}}

\usepackage{hyperref}
\hypersetup{colorlinks=false, pdfborder={0 0 0}}

% Настройка маркированных списков (тире с висячим отступом)
\usepackage{enumitem}
% \usepackage{enumitem}

\setlist[itemize]{
    label=---,
    leftmargin=1.25pt,            % Общий отступ списка = 0
    itemindent=1.92cm,         % Первая строка с красной строки
    listparindent=0pt,         % Вторая строка без отступа
    % labelsep=0.5em,
    itemsep=0cm,
    topsep=0pt,
    partopsep=0pt,
    parsep=0cm,
    after=\vspace{0em},
    % label=---,
    wide=1.25cm,               % Начало маркера строго с красной строки
    % leftmargin=1.25cm,         % Вторая и последующие строки выровнены по маркеру
    % labelsep=0.3em,            % Небольшой пробел между тире и текстом
    % itemsep=0pt,
    % topsep=0pt,
    % partopsep=0pt,
    % parsep=0pt
}

% Настройка нумерованных списков (enumerate) - если нужно тоже с тире
\setlist[enumerate]{
    label=(\arabic*),
    leftmargin=1.25pt,
    itemindent=1.92cm,
    listparindent=0pt,
    labelsep=0.5em,
    itemsep=0pt,
    topsep=0pt,
    partopsep=0pt,
    parsep=0pt,
    after=\vspace{0em}, % ← Убирает отступ после списка
    label=\arabic*.,
    wide=1.25cm,               % Начало цифры строго с красной строки
    % leftmargin=1.25cm,         % Вторая строка выровнена ровно под цифру
    % labelsep=0.4em,            % Пробел после скобки
    % itemsep=0pt,
    % topsep=0pt,
    % partopsep=0pt,
    % parsep=0pt
}


% НАСТРОЙКА РИСУНКОВ
\usepackage[russian]{babel}
\usepackage{caption}

% Переопределение формата: Рисунок 1 - Название
\DeclareCaptionLabelSeparator{hyphen}{ --- }
\captionsetup[figure]{labelsep=hyphen, name=Рисунок, justification=centering,}

% \usepackage{caption}

% % Настройка подписей к рисункам
% \captionsetup[figure]{
%     name=Рисунок,
%     labelsep=period,
%     font=normal,
%     skip=1em,
%     justification=centering,
%     singlelinecheck=off
% }
% \captionsetup[figure]{
%     name=Рисунок,
%     labelsep=hyphen,
% }

% ====== ПРИМЕР ТАБЛИЦЫ ======
% \begin{table}[htbp]
% \centering
% \caption{Деградация метрик качества при увеличении уровня шума}
% \begin{tabularx}{\textwidth}{|c|c|c|c|c|>{\centering\arraybackslash}X|}
% \hline
% Уровень шума & RMSE & NASA Score & F1-мера & AUC-ROC & Относительное падение F1 \\
% \hline
% 0\% & 0.031 & 287 & 0.89 & 0.94 & - \\
% 2\% & 0.034 & 301 & 0.87 & 0.92 & 2.2\% \\
% 5\% & 0.041 & 338 & 0.83 & 0.88 & 6.7\% \\
% 10\% & 0.056 & 402 & 0.74 & 0.79 & 16.8\% \\
% \hline
% \end{tabularx}
% \label{tab:noise_robustness}
% \end{table}

% ==========================================
% НАЧАЛО ДОКУМЕНТА
% ==========================================

\begin{document}
% Убираем нумерацию для содержания
\makeatletter
\renewcommand{\numberline}[1]{}
\makeatother


% ==========================================
% НУМЕРАЦИЯ НАЧИНАЕТСЯ С 6-Й СТРАНИЦЫ
% ==========================================
\pagenumbering{arabic}
\setcounter{page}{6}

\tableofcontents
\clearpage



% ==========================================
% ВВЕДЕНИЕ
% ==========================================
\section{ВВЕДЕНИЕ}
\vspace{1em}
В современных условиях эксплуатации сложных технологических систем крайне важно не только учитывать текущее состояние составляющих систему модулей, но и уметь прогнозировать их поведение на несколько эксплуатационных циклов вперёд. Прогнозирование выхода параметров за допустимые границы позволяет не только снизить риски аварийной эксплуатации, но и повысить общую эффективность производственного процесса, предотвратив экономические издержки, связанные с незапланированными простоями или преждевременным снятием с эксплуатации ещё ресурсных узлов.

Оценка и моделирование текущего состояния оборудования являются ключевыми факторами перехода к предиктивному обслуживанию. Своевременное формирование вероятностных сценариев развития состояния дает возможность использовать оборудование с максимальной эффективностью и оптимизировать графики технического обслуживания. В работах [2–10] рассматривается решение проблем обнаружения аномалий и прогнозирования остаточного ресурса технологического оборудования с использованием моделей глубокого обучения. Особое внимание в современных исследованиях уделяется генеративным архитектурам, способным моделировать распределения временных рядов и синтезировать правдоподобные траектории развития отказов.

Особенностью данной работы является применение вариационного автокодировщика (VAE) в качестве основы системы. Модель решает две сопряженные задачи: классическую задачу восстановления текущих параметров и задачу генерации последующих n эксплуатационных циклов. Амортизация параметров латентного распределения z выполняется не по полному временному ряду, а на основе начальных n циклов, что позволяет эффективно оперировать скалярными показателями состояния St (не являющимися векторами) и генерировать элементы выходной последовательности. 
Математический аппарат алгоритма базируется на адаптации вариационной нижней оценки (ELBO) под задачу последовательного прогнозирования.

Данная работа имеет практическую значимость в областях производства и эксплуатации технологически сложного оборудования, в частности турбовентиляторных двигателей и промышленных агрегатов. Её применение позволяет формировать прогнозы будущих показателей циклов, выявлять скрытые закономерности деградации и поддерживать принятие решений в условиях неполных или зашумленных данных.

В связи с этим целью данной работы являлось снижение эксплуатационных рисков и повышение экономической эффективности за счет создания системы моделирования состояния оборудования на базе вариационного автокодировщика, обеспечивающей генерацию последующих циклов и точную оценку технического состояния.

\clearpage



% ==========================================
% РАЗДЕЛ I
% ==========================================
\section{1 АНАЛИЗ ПРЕДМЕТНОЙ ОБЛАСТИ И СУЩЕСТВУЮЩИХ АНАЛОГОВ}
\section{1.1 Обзор исследований и моделей машинного обучения для решения аналогичных задач}
\section{1.1.1 Применение классических и глубоких автокодировщиков для оценки состояния оборудования}
\vspace{1em}
В работе Jia Guo, Guannan Liu, Yuan Zuo, Junjie Wu под названием “An Anomaly Detection Framework Based on Autoencoder and Nearest Neighbor” [2] исследователи предлагают архитектуру AEKNN, объединяющую автокодировщик и метод k-ближайших соседей. На этапе обучения модель тренируется исключительно на данных штатного режима. Сжатое представление в латентном слое используется как вход для k-NN классификатора. При тестировании аномалии выявляются по отклонению расстояния до ближайших соседей в скрытом пространстве. Эффективность подхода подтверждается на промышленных датасетах, однако метод не предусматривает генерацию будущих состояний и работает в режиме классификации «штатный/аварийный» без оценки вероятностной неопределенности.

В статье «Обнаружение аномальных событий на хосте с использованием автокодировщика» авторы Гурина А. О., Гузев О. Ю., Елисеев В. Л. предлагают метод построения модели нормального поведения системы на основе двух параллельных автокодировщиков с различной архитектурой. Критерием аномальности выступает превышение порога мгновенной ошибки реконструкции (IRE), рассчитываемого отдельно для каждого декодера. Метод демонстрирует высокую чувствительность к редким событиям, но опирается на детерминированную реконструкцию, что ограничивает его применимость для стохастических процессов износа технологического оборудования.


\vspace{1em}
\section{1.1.2 Вариационные автокодировщики в задачах вероятностного прогнозирования и генерации временных рядов}
\vspace{1em}
В работе Kingma and Welling “Auto-Encoding Variational Bayes” [16] заложены теоретические основы VAE, где максимизация вариационной нижней оценки (ELBO) обеспечивает одновременное обучение кодировщика, декодировщика и регуляризацию латентного пространства. Применительно к телеметрии оборудования данная архитектура позволяет оценивать не только точечное восстановление параметров, но и дисперсию реконструкции, что напрямую связано с надежностью прогноза.

Исследование Chung et al. “Deep Generative Models for Time Series Forecasting” демонстрирует применение рекуррентного VAE (VRNN) для многошагового прогнозирования. Авторы показывают, что аморизованный вывод параметров распределения $z$ по скользящему окну данных позволяет генерировать правдоподобные траектории развития параметров. Однако в работе используется векторное представление состояния на каждом шаге, что увеличивает вычислительную нагрузку и усложняет интерпретацию скалярных индикаторов работоспособности

\vspace{1em}
\section{1.1.3 Сравнительный анализ генеративных архитектур (GAN, Diffusion, VAE) для моделирования циклических процессов}
\vspace{1em}
В обзоре Goodfellow et al. по генеративным состязательным сетям и последующих работах по Diffusion-моделям отмечается высокое качество генерации сложных распределений. Однако для задач прогнозирования циклов технологического оборудования эти архитектуры демонстрируют существенные ограничения: GAN страдают от нестабильности обучения (mode collapse) и отсутствия явной вероятностной интерпретации латентных переменных, что затрудняет калибровку порогов аварийных состояний. Diffusion-модели требуют многоэтапного обратного процесса генерации, что критично увеличивает время инференса и делает их неоптимальными для систем мониторинга в реальном времени.
В работе Sohn et al. “Learning Structured Output Representation using Deep Conditional Generative Models” [20] показано, что условные VAE (CVAE) обеспечивают стабильную оптимизацию и явную связь между входным контекстом (текущие n циклов) и выходной последовательностью. Это позволяет формировать прогнозы с оценкой доверительных интервалов, что напрямую соответствует требованиям к системам предиктивного обслуживания.

\vspace{1em}
\section{1.1.4 Адаптация функции ELBO и амортизированный вывод для многошагового прогнозирования состояний}
\vspace{1em}
В исследованиях по байесовским методам прогнозирования [21, 22] отмечается, что стандартная формулировка ELBO требует модификации для последовательных данных. Авторы предлагают аморизовывать параметры распределения не по полному временному ряду, а по начальному окну из nциклов, что снижает риск переобучения и позволяет декодировщику генерировать последующие элементы на основе фиксированного латентного представления.
В работе Zheng et al. “Variational Autoencoder for Remaining Useful Life Prediction” скалярный показатель остаточного ресурса выводится как функция от дисперсии реконструкции и расстояния в латентном пространстве. Авторы подчеркивают, что не является вектором признаков, а представляет собой агрегированную метрику, вычисляемую на основе вероятностного вывода. Такой подход позволяет отделить задачу генерации телеметрии от задачи оценки состояния, сохраняя при этом их математическую связность через общий ELBO.

\vspace{1em}
\section{1.2 Критерии оценки и валидации генеративных моделей в промышленной диагностике}
\vspace{1em}
Валидация генеративных моделей для временных рядов оборудования требует сочетания поточечных и структурных метрик. В работах [24, 25] базовой метрикой выступает поточечная MSE (Mean Squared Error), вычисляемая для каждого шага прогнозирования: размерность вектора наблюдений. Для учёта фазовых сдвигов и нелинейных деформаций временных осей дополнительно применяются метрики DTW (Dynamic Time Warping) и MAE.
В исследовании показано, что для оценки качества генерации будущих циклов недостаточно лишь минимизации MSE, так как модель может «усреднять» распределение. Поэтому вводится критерий коэффициента детерминации по каждому прогнозируемому циклу и F1-score классификации режимов, где порог аномальности определяется на основе распределения ошибок реконструкции и расстояний Махаланобиса в латентном пространстве. Комплексное применение данных метрик позволяет объективно оценить как точность восстановления, так и устойчивость генерации к зашумленным телеметрическим данным.

\vspace{1em}
\section{1.3 Научная новизна}
\vspace{1em}
% \addcontentsline{toc}{subsection}{Научная новизна}
\begin{itemize}
    \item Разработана оригинальная модификация конвейера вариационного автоэнкодера (TimeSeries Denoising VAE), совмещающая в себе принципы шумоподавления (Denoising) на этапе кодирования предыстории и пошаговый авторегрессионный (Closed-loop) механизм генерации многомерного будущего. В отличие от классических VAE, предсказывающих весь горизонт прогнозирования единой матрицей и приводящих к математическому коллапсу (усреднению траекторий до константных прямых), предложенный подход позволяет рекурсивно моделировать нелинейную физику деградации датчиков. 
    \item Предложена стохастическая схема пошагового сэмплирования динамических априорных распределений (\(Z_{t}\)) в скрытом пространстве рекуррентных и марковских моделей (VRNN, Deep SSM), адаптированная для оценки остаточного ресурса (RUL) авиационных двигателей в условиях неопределенности. Это позволило воссоздавать не просто детерминированный тренд износа, а реалистичные физические колебания параметров (микро-флуктуации датчиков), отражающие реальные условия эксплуатации.
\end{itemize}

\vspace{1em}
\section{1.4 Теоретическая и практическая значимость}
\vspace{1em}
Теоретическая значимость заключается в обосновании применения вариационного автокодировщика для совместного решения задач многошагового прогнозирования и вероятностной оценки надежности. За счет совместной оптимизации точности генерации будущих циклов и регуляризации скрытого пространства дивергенцией Кульбака-Лейблера обеспечен переход от детерминированного прогнозирования к стохастической оценке состояния, позволяющей количественно измерять неопределенность и риск отказа в условиях отсутствия данных об аварийных режимах.»
Полученные результаты создают теоретическую основу для построения систем предиктивного обслуживания, способных количественно оценивать неопределенность прогноза и адаптироваться к изменяющимся условиям эксплуатации без необходимости переобучения на аварийных сценариях.

Практическая значимость работы заключается в разработке открытого программного комплекса в виде Python-библиотеки, реализующего полный жизненный цикл системы предиктивного обслуживания оборудования. Разработанный инструмент обеспечивает следующее:
\begin{enumerate}
    \item Обеспечение масштабируемости обработки телеметрических данных. За счет использования генераторов при чтении данных реализована конвейерная обработка потоков информации, что позволяет работать с большими массивами телеметрии без полной выгрузки их в оперативную память. Это критически важно для задач мониторинга парков оборудования, где объемы данных измеряются терабайтами.
    \item Ускорение внедрения методов машинного обучения. Модульная структура библиотеки, включающая пакеты предобработки и автоматизированные пайплайны, снижает порог входа для инженеров-разработчиков. Стандартизированные процедуры очистки и нормализации данных минимизируют количество ручных операций при подготовке датасетов.
    \item Промышленная готовность и воспроизводимость. Интеграция сервисов MLflow позволяет осуществлять версионирование моделей, логирование гиперпараметров экспериментов и централизованное хранение артефактов. Это обеспечивает прозрачность процесса обучения и возможность быстрого отката к предыдущим версиям моделей при изменении условий эксплуатации.
    \item Упрощение процесса принятия решений. В состав библиотеки включены специализированные методы (например, автоматический расчет кумулятивной функции распределения ошибок и прогнозирование остаточного ресурса), которые абстрагируют сложные математические вычисления. Это позволяет эксплуатационному персоналу получать интерпретируемые метрики риска отказа без необходимости глубокого понимания внутреннего устройства нейросетевых архитектур.
    \item Возможность независимого использования модулей. Архитектура библиотеки спроектирована таким образом, что каждый пакет (чтение данных, метрики, модели) может быть использован автономно или интегрирован в существующие корпоративные системы мониторинга через программный интерфейс (API).
\end{enumerate}
Валидация разработанных алгоритмов и программной реализации выполнена на базе открытого набора данных NASA C-MAPSS. Полученные результаты подтверждают корректность работы модулей прогнозирования и создают основу для переноса методики на задачи мониторинга реального оборудования.

\vspace{1em}
\section{1.5 Положения, выносимые на защиту}
\vspace{1em}
% Использование вариационного автокодировщика, обученного на данных нормального функционирования, в сочетании с оценкой эмпирической функции распределения ошибок реконструкции позволяет эффективно детектировать отклонения в работе оборудования и прогнозировать остаточный ресурс с заданной вероятностью, превосходя по устойчивости к шуму методы детерминированного прогнозирования в условиях отсутствия разметки данных об отказах.
Метод предиктивного моделирования деградации сложных технических систем на базе оригинальной модификации вариационного автоэнкодера (TimeSeries Denoising VAE).
В отличие от традиционных вероятностных подходов прямого прогнозирования, приводящих к математическому коллапсу дисперсии и аппроксимации будущего плоскими константными линиями, предложенный метод реализует концепцию пошагового авторегрессионного инференса с глубокой обратной связью (Closed-loop forecasting).
Это обеспечивает устойчивое воссоздание долгосрочных, сложновыраженных нелинейных траекторий износа многомерных временных рядов (показаний датчиков) с сохранением их физических законов затухания и роста.
\vspace{1em}
\section{1.6 Выводы}
\vspace{1em}

В рамках первой главы выполнен анализ предметной области прогнозирования технического состояния сложного оборудования и проведен сравнительный обзор существующих методов машинного обучения. На основе проведенного исследования сформулированы следующие основные результаты:

\begin{enumerate}
\item Установлено, что традиционные детерминированные методы прогнозирования не обеспечивают оценки неопределенности прогноза, что критично для задач мониторинга ответственных узлов. Генеративные состязательные сети и диффузионные модели, несмотря на высокую точность генерации, обладают значительной вычислительной сложностью, затрудняющей их применение в системах реального времени.

\item Предложена концепция многошагового прогнозирования на основе амортизированного вывода параметров скрытого распределения. Использование скользящего окна из $n$ предыдущих циклов позволяет модели агрегировать контекст временного ряда, снижая влияние локальных шумов телеметрии на точность оценки текущего состояния.

\item Разработан методологический подход к детектированию аномалий без использования размеченных данных об отказах. Суть подхода заключается в обучении модели исключительно на данных нормального функционирования и последующем вычислении вероятности отказа через эмпирическую функцию распределения ошибок реконструкции.

\item Сформулированы функциональные и нефункциональные требования к программной реализации системы. Определена модульная архитектура комплекса, включающая конвейер обработки данных, блок генеративного моделирования и сервисы управления экспериментами, что обеспечивает масштабируемость и воспроизводимость результатов.
\end{enumerate}

Полученные теоретические и методологические результаты создают необходимую базу для перехода к практической реализации. Во второй главе будет выполнено детальное проектирование системы моделирования, формализована математическая постановка задачи и разработана архитектура программного комплекса.

\clearpage



% ==========================================
% РАЗДЕЛ II
% ==========================================
\section{2 ПРОЕКТИРОВАНИЕ СИСТЕМЫ МОДЕЛИРОВАНИЯ ДАННЫХ}
\section{2.1 Постановка задачи}
\vspace{1em}

Для создания системы предиктивного моделирования состояния технологического оборудования в качестве базового источника данных принят датасет NASA C-MAPSS (Commercial Modular Aero-Propulsion System Simulation), содержащий телеметрические показания турбовентиляторных двигателей в различных режимах эксплуатации и условиях окружающей среды.

Требуется спроектировать и реализовать программный комплекс для вероятностного моделирования и прогнозирования состояния оборудования на основе нейросетевых генеративных моделей. Основными требованиями, предъявляемыми к системе, являются:

\begin{enumerate}
\item Система должна осуществлять многошаговое прогнозирование телеметрических параметров на $n$ эксплуатационных циклов вперёд с оценкой вероятностной неопределённости.
\item Программная реализация должна базироваться на архитектуре вариационного автокодировщика (VAE), обеспечивающей явное вероятностное латентное пространство и устойчивую оптимизацию.
\item Система должна принимать на вход временные ряды показаний датчиков в форматах CSV/Parquet. Единицей обработки выступает скользящее окно из $n$ последовательных циклов $X_{1:n}$.
\item Система должна формировать скалярный индикатор состояния $S_t$ на основе статистик латентного пространства и ошибки реконструкции, исключая векторное представление статуса оборудования.
\item Модель должна поддерживать модульную интеграцию: загрузку предобученных весов из удалённого хранилища и экспорт предсказаний в форматы, совместимые с промышленными SCADA/MES-системами.
\item Все программные компоненты должны быть оформлены в виде независимых Python-пакетов с чётким разделением ответственности (предобработка, обучение, инференс, валидация).
\item В системе должна быть реализована автоматизированная регистрация метрик обучения, гиперпараметров и версий датасетов для обеспечения воспроизводимости экспериментов.
\item Оценка качества модели должна производиться по комплексу метрик: поточечная MSE, DTW для временных рядов, а также классификационные показатели на основе порогов $S_t$.
\end{enumerate}

\section{2.1.1 Описание входных данных}
Исходным информационным базисом для проведения исследований является плоский двумерный массив полетной телеметрии, содержащий последовательные записи о состоянии парка авиационных двигателей NASA Turbofans (C-MAPSS). Вектор сырых данных для каждого отдельного двигателя представляет собой временную последовательность полетных циклов.

Пусть $\mathbf{X}_{\text{raw}} \in \mathbb{R}^{N_{\text{total}} \times (2 + D)}$ — исходная плоская матрица данных, где $N_{\text{total}} = 86914$ — общее число зарегистрированных полетных циклов (строк) в выборке, а размерность столбцов включает в себя два служебных идентификатора и $D = 26$ информационных каналов (операционные настройки и показания физических датчиков). Спецификация столбцов исходной матрицы имеет вид:
\begin{equation}
\mathbf{X}_{\text{raw}} = \left[ \mathbf{u}, \mathbf{c}, \mathbf{f}_1, \mathbf{f}_2, \dots, \mathbf{f}_D \right],
\end{equation}
где $\mathbf{u} \in \mathbb{N}^{N_{\text{total}}}$ - вектор идентификаторов контролируемых двигателей (\textit{unit number}), $\mathbf{c} \in \mathbb{N}^{N_{\text{total}}}$ - вектор порядковых номеров полетных циклов (\textit{time in cycles}), а $\mathbf{f}_d \in \mathbb{R}^{N_{\text{total}}}$ - вектор физических измерений $d$-го информационного канала.

В рамках разработанного конвейера предобработки служебные столбцы ($\mathbf{u}, \mathbf{c}$) изолируются от процедуры внесения искажений, а матрица датчиков подвергается Z-нормализации (\textit{StandardScaler}) и процедуре скользящего сглаживания (\textit{Rolling Mean}) с размером окна в 5 циклов. Это позволяет сформировать эталонную плоскую матрицу чистых физических трендов:
\begin{equation}
\mathbf{X}_{\text{clean}} = \text{StandardScaler}\left( \text{RollingMean}\left( \left[ \mathbf{f}_1, \dots, \mathbf{f}_D \right] \right) \right) \in \mathbb{R}^{N_{\text{total}} \times D}.
\end{equation}

Для проведения стресс-тестирования моделей на устойчивость к отказам бортового оборудования формируется искаженная матрица входных признаков $\mathbf{X}_{\text{stressed}} = \text{apply\_stress\_to\_data}(\mathbf{X}_{\text{clean}})$. На основе переданного параметра \textit{noise\_type} на отмасштабированные датчики накладываются два типа аппаратных дефектов:

% \textit{1. Высокочастотный белый шум (\textit{noise\_type = 'noise'}).} \\
К каждому элементу матрицы чистых признаков добавляется случайная составляющая, подчиняющаяся нормальному (Гауссову) закону распределения с нулевым математическим ожиданием и заданным уровнем стандартного отклонения $\sigma_{\text{noise}} = 0.1$:
\begin{equation}
x_{\text{stressed}, i, d} = x_{\text{clean}, i, d} + \epsilon_{i, d}, \quad \epsilon_{i, d} \sim \mathcal{N}(0, \sigma_{\text{noise}}^2).
\end{equation}

% \textit{2. Аппаратные пропуски и обрывы связи (\textit{noise\_type = 'drop'}).} \\
Симулируется случайная потеря пакетов телеметрии с интенсивностью $\text{drop\_rate} = 0.2$. Ровно $20\,\%$ случайных ячеек матрицы датчиков заменяются на индикатор полного отсутствия сигнала $\text{fill\_value} = 0.0$:
\begin{equation}
x_{\text{stressed}, i, d} = \begin{cases} 
0.0, & \text{с вероятностью } 0.2, \\
x_{\text{clean}, i, d}, & \text{с вероятностью } 0.8.
\end{cases}
\end{equation}

В комбинированном режиме оба дефекта накладываются последовательно, после чего значения принудительно ограничиваются методом \textit{clip} для предотвращения математических выбросов.

На финальном этапе плоские матрицы нарезаются на скользящие трехмерные окна с заданной глубиной и значением сдвига. В результате формируются две ключевые структуры данных, поступающие на вход нейросети:
\begin{enumerate}
    \item Искаженная предыстория (Вход Энкодера): трехмерный тензор $\mathcal{X}_{\text{stressed}} \in \mathbb{R}^{B \times M \times D}$, содержащий зашумленные и рваные сигналы первых 10 циклов скользящего окна. На основе этой матрицы энкодер учится строить инвариантный латентный код стиля.
    \item Идеальная предыстория (Точка опоры Декодера): изолированный вектор $\mathbf{x}_{\text{anchor}} \in \mathbb{R}^{B \times D}$, представляющий собой строго крайнее, отмасштабированное значение из сглаженной (чистой) матрицы $\mathbf{X}_{\text{clean}}$. Данный вектор служит безошибочным физическим якорем, от которого декодер рекурсивно вычисляет приращения дельт в будущее.
\end{enumerate}

Здесь $B = 87894$ — результирующее количество сформированных трехмерных окон (размерность батча) после фильтрации граничных циклов индивидуальных двигателей.

\section{2.1.2 Описание выходных данных}

Целевым результатом работы системы является генерация вероятностного прогноза будущего состояния оборудования. Выход модели представляет собой последовательность сгенерированных векторов наблюдений для $k$ будущих циклов:

Пусть $D = 26$ — размерность пространства измеряемых признаков (каналов датчиков), $M = 10$ — глубина анализируемой предыстории (длина входного зашумленного окна), а $T = 20$ — полный временной горизонт, включающий фазу восстановления и фазу прогнозирования.

% Матричная форма \(\mathbb{R}^{T\times D}\) показывает строгое многомерное описание выходных параметров.Формула рекуррентного перехода через \(\mathbf{\Delta }_{t}^{(s)}\) математически описывает, почему модель генерирует плавно изгибающийся физический тренд, а не падает в константное усреднение плоских линий [INDEX, INDEX].Введение операторов матожидания и дисперсии (\(\={\mathbf{y}}_{t}\) и \(\mathbf{\sigma }_{t}^{2}\)) переводит задачу из плоскости «угадывания чисел» в плоскость строгой интервальной оценки рисков отказа в предиктивном обслуживании.

Каждый $s$-й сгенерированный моделью сценарий представляет собой многомерный временной ряд, описываемый матрицей $\hat{\mathbf{Y}}^{(s)}$ размерности $(T \times D)$:
\begin{equation}
\hat{\mathbf{Y}}^{(s)} = \begin{bmatrix} 
\hat{y}_{1,1}^{(s)} & \hat{y}_{1,2}^{(s)} & \dots & \hat{y}_{1,D}^{(s)} \\
\hat{y}_{2,1}^{(s)} & \hat{y}_{2,2}^{(s)} & \dots & \hat{y}_{2,D}^{(s)} \\
\vdots & \vdots & \ddots & \vdots \\
\hat{y}_{T,1}^{(s)} & \hat{y}_{T,2}^{(s)} & \dots & \hat{y}_{T,D}^{(s)}
\end{bmatrix} \in \mathbb{R}^{T \times D}, \quad s \in \{1, 2, \dots, S\},
\end{equation}
где $S = 5$ — общее количество генерируемых альтернативных сценариев для оценки интервальной неопределенности, а $\hat{y}_{t,d}^{(s)}$ — нормализованное (StandardScaler) значение $d$-го датчика в момент времени $t$ в рамках $s$-го сценария.

Структурно выходная матрица $\hat{\mathbf{Y}}^{(s)}$ разделена на две функциональные зоны вдоль временной оси $t$:
\begin{enumerate}
    \item {Зона деноизинга и реконструкции ($1 \le t \le M$):} фаза, на которой декодер восстанавливает истинные значения на основе латентного вектора стиля $\mathbf{z}^{(s)}$. В силу использования Denoising-подхода, выходные значения $\hat{y}_{t,d}^{(s)}$ аппроксимируют сглаженный (Rolling Mean) физический тренд деградации $\mathbf{Y}_{\text{smooth}}$, полностью игнорируя наложенный на входные данные $\mathbf{X}_{\text{stressed}}$ измерительный хаос и пропуски:
    \begin{equation}
    \hat{\mathbf{Y}}_{1:M}^{(s)} \approx \mathbf{Y}_{\text{smooth}, 1:M}, \quad \text{при этом} \quad \hat{\mathbf{Y}}_{1:M}^{(1)} \equiv \hat{\mathbf{Y}}_{1:M}^{(2)} \equiv \dots \equiv \hat{\mathbf{Y}}_{1:M}^{(S)}.
    \end{equation}
    Траектории всех сценариев в этой зоне идентичны, так как они жестко привязаны к детерминированной предыстории конкретного двигателя.
    
    \item {Зона авторегрессионного прогноза ($M < t \le T$):} фаза автономной генерации будущего с глубокой обратной связью (Closed-loop). Значение на шаге $t$ формируется рекурсивно на основе вектора приращений (дельт) $\mathbf{\Delta}_t^{(s)}$, вычисляемого моделью на базе накопленной предыстории скрытых состояний:
    \begin{equation}
    \hat{\mathbf{y}}_{t}^{(s)} = \hat{\mathbf{y}}_{t-1}^{(s)} + \mathbf{\Delta}_t^{(s)}\left(\hat{\mathbf{y}}_{1:t-1}^{(s)}, \mathbf{z}_t^{(s)}\right), \quad t \in \{M+1, \dots, T\}.
    \end{equation}
\end{enumerate}

Для получения финального точечного прогноза и оценки дисперсии (интервальной неопределенности) пооконных траекторий на диссертационную защиту выносится использование операторов математического ожидания $\mathbb{E}$ и дисперсии $\mathbb{D}$ над множеством сценариев $S$:
\begin{equation}
\bar{\mathbf{y}}_{t} = \mathbb{E}_{s} \left[ \hat{\mathbf{y}}_{t}^{(s)} \right] = \frac{1}{S} \sum_{s=1}^{S} \hat{\mathbf{y}}_{t}^{(s)},
\end{equation}
\begin{equation}
\mathbf{\sigma}_{t}^2 = \mathbb{D}_{s} \left[ \hat{\mathbf{y}}_{t}^{(s)} \right] = \frac{1}{S-1} \sum_{s=1}^{S} \left( \hat{\mathbf{y}}_{t}^{(s)} - \bar{\mathbf{y}}_{t} \right)^2.
\end{equation}

Величина $\mathbf{\sigma}_{t}^2 \in \mathbb{R}^D$ математически описывает расхождение веера траекторий. Рост дисперсии во времени ($\sigma_{t_2}^2 > \sigma_{t_1}^2$ при $t_2 > t_1 > M$) отражает экспоненциальное накопление неопределенности прогноза остаточного ресурса (RUL) по мере удаления от последней известной физической точки отсчета технического состояния авиадвигателя.





\end{document}